\subsection{APPLICATION TO LINEAR REPRESENTATIONS}

PROPOSITION 13. Let G be a connected Lie group and \(\rho\) an analytic linear representation of G on a complete normable space E. Let \(\mathrm{E}_1, \mathrm{E}_2\) be two closed vector subspaces of E such that \(\mathrm{E}_2 \subset \mathrm{E}_1\) . The following conditions are equivalent:

\begin{enumerate}
\def\labelenumi{(\roman{enumi})}
\item
  \(\rho(g)x \equiv x \pmod{\mathbf{E}_2}\) for all \(g \in \mathbf{G}\) and all \(x \in \mathbf{E}_1\) ;
\item
  \(\mathrm{L}(\rho)(\mathrm{L}(\mathrm{G}))\) maps \(E_{1}\) into \(E_{2}\) .
\end{enumerate}

\[
\begin{array}{l l l} \rho (g) x \equiv x & (\mathrm{mod} \mathrm{E} _ {2}) & \text {for all \(g\in G\) and all \(x\in E_{1}\)} \\ \Leftrightarrow \rho (\exp a) x \equiv x & (\mathrm{mod} \mathrm{E} _ {2}) & \text {for all \(a\in L(G)\) and all \(x\in E_{1}\)} \\ \Leftrightarrow (\exp \mathrm{L} (\rho) a) x \equiv x & (\mathrm{mod} \mathrm{E} _ {2}) & \text {for all \(a\in L(G)\) and all \(x\in E_{1}\)}. \end{array}
\]

On the other hand, if \(u \in \mathcal{L}(\mathrm{E})\) , then

\[
\begin{array}{r l} & \exp (\lambda u) x \equiv x \pmod {\mathrm{E} _ {2}} \quad \text {(for all} \lambda \in \mathrm{K} \text {and all} x \in \mathrm{E} _ {1} \\ & \Leftrightarrow u (\mathrm{E} _ {1}) \subset \mathrm{E} _ {2} \end{array}
\]

whence the proposition.

COROLLARY 1. For \(\mathbf{E}_1\) to be stable under \(\rho\) , it is necessary and sufficient that \(\mathbf{E}_1\) be stable under \(\mathrm{L}(\rho)\) .

It suffices to take \(E_{1} = E_{2}\) in Proposition 13.

COROLLARY 2. Suppose that \(\rho\) is finite-dimensional. For \(\rho\) to be simple (resp. semi-simple), it is necessary and sufficient that \(\mathbf{L}(\rho)\) be simple (resp. semi-simple).

This follows from Corollary 1.

COROLLARY 3. Let \(x \in E\) . For \(x\) to be invariant under \(\rho(G)\) , it is necessary and sufficient that \(x\) be annihilated by \(L(\rho)(L(G))\) (that is that \(x\) be invariant under \(L(\rho)\) in the sense of Chapter I, § 3, Definition 3).

It suffices to take \(\mathbf{E}_1 = \mathbf{K}x\) and \(\mathbf{E}_2 = 0\) in Proposition 13.

COROLLARY 4. Let \(\rho'\) be another analytic linear representation of \(G\) on a complete normable space \(E'\) . Let \(T \in \mathcal{L}(E, E')\) . The following conditions are equivalent:

\begin{enumerate}
\def\labelenumi{(\roman{enumi})}
\item
  \(\mathrm{T}\rho (g) = \rho '(g)\mathrm{T}\) for all \(g\in \mathbf{G}\)
\item
  \(\mathrm{TL}(\rho)(a) = \mathrm{L}(\rho')(a)\mathrm{T}\) for all \(a\in \mathrm{L}(\mathrm{G})\)
\end{enumerate}

Let \(\sigma\) be the linear representation of \(G\) on \(\mathcal{L}(E, E')\) derived from \(\rho\) and \(\rho'\) ( \(\S 3\) , no. 11, Corollary 1 to Proposition 41). Condition (i) means that \(T\) is invariant under \(\sigma(G)\) . Condition (ii) means that \(T\) is annihilated by \(L(\sigma)(L(G))\) . It then suffices to apply Corollary 3.

COROLLARY 5. Suppose that \(\rho\) and \(\rho'\) are finite-dimensional. For \(\rho\) and \(\rho'\) to be equivalent, it is necessary and sufficient that \(\mathbf{L}(\rho)\) and \(\mathbf{L}(\rho')\) be equivalent.

This is a special case of Corollary 4.

COROLLARY 6. Suppose that \(\mathbf{G}\) is finite-dimensional. Let \(t \in \mathrm{U}(\mathrm{G})\) . For \(\mathbf{L}_t\) (resp. \(\mathbf{R}_t\) ) to be right (resp. left) invariant, it is necessary and sufficient that \(t\) belong to the centre of \(\mathrm{U}(\mathrm{G})\) .

For \(\mathrm{L}_t\) (resp. \(\mathbf{R}_t\) ) to be right (resp. left) invariant, it is necessary and sufficient that \(\varepsilon_g * t = t * \varepsilon_g\) for all \(g \in G\) , that is that \((\operatorname{Int} g)_* t = t\) . There exists an integer \(n\) such that \(t \in \mathrm{U}_n(G)\) . By Corollary 3 and Proposition 45 of § 3, no. 12, \((\operatorname{Int} g)_* t = t\) for all \(g \in G\) if and only if \([a, t] = 0\) for all \(a \in \mathrm{L}(G)\) , that is if and only if \(t\) commutes with \(\mathrm{U}(G)\) .

